\chapter{The orbit quotient and the propositional encoding}\label{sec:encoding}

By Theorem~\ref{thm:reduction} it suffices to decide whether some graph
invariant under the canonical permutation $\sigma_0$ is strongly regular with
parameters $(99,14,1,2)$.  This chapter carries out the reduction of that
question to a single propositional formula $\Phi$, and proves the direction
that the certificate method needs: every invariant srg would give rise to a
satisfying valuation of (the pre-normalisation part of) $\Phi$.  The chapter
has five strands: the orbit quotient of the edge set (\S\ref{ssec:enc-orbits}),
the formula itself as a mathematical object, a Hoare-style framework for
reasoning about clause-emitting transformations, the soundness of the three
counting gadgets, and finally the bridge from strong regularity to the
semantic counting invariants.

\section{The orbit quotient}\label{ssec:enc-orbits}

\begin{definition}[Pair orbits and representatives]\label{def:enc-orbitmap}
The cyclic group $\langle\sigma_0\rangle$ acts on unordered pairs of vertices
$\{u,v\} \subseteq \{0,\dots,98\}$, $u \neq v$, by
$\sigma_0\cdot\{u,v\} = \{\sigma_0(u),\sigma_0(v)\}$.  Write
$\mathrm{orb}(u,v)$ for the index of the orbit of $\{u,v\}$, orbits being
numbered $0,1,2,\dots$ in order of first appearance as the pairs are
enumerated lexicographically, and for each orbit index $o$ let
$\mathrm{rep}(o)$ denote the first pair of orbit $o$ in that enumeration (its
\emph{representative}).  The map $\mathrm{orb}$ is symmetric in its two
arguments and constant along the action:
$\mathrm{orb}(\sigma_0 u, \sigma_0 v) = \mathrm{orb}(u,v)$.
\lean{ConwayO7.orbitOf, ConwayO7.orbitRep, ConwayO7.orbitOf_swap,
ConwayO7.orbitOf_sigma0}
\leanfile{ConwayO7/Orbits.lean}\leanok
\uses{def:sigma0}
\end{definition}

\begin{lemma}[Orbit structure of pairs]\label{lem:orbits}
The action of $\langle\sigma_0\rangle \cong \ZZ/7$ on the $\binom{99}{2}
= 4851$ unordered pairs of vertices has exactly $693$ orbits, each of size
$7$; equivalently, $\mathrm{orb}$ takes exactly the values $0,\dots,692$, and
every pair is reached from its orbit's representative by iterating the
action.
\lean{ConwayO7.numOrbits_eq, ConwayO7.orbitOf_lt,
ConwayO7.orbitRep_reaches}
\leanfile{ConwayO7/Orbits.lean}\leanok
\uses{def:enc-orbitmap}
\end{lemma}

\begin{proof}\leanok
No pair is fixed by $\sigma_0$ (its two $7$-cycles share no point with any
pair twice), so every orbit has size exactly $7$ and $4851/7 = 693$.  In the
formalisation the orbit table is computed once and the three statements are
verified by evaluation.
\end{proof}

\begin{definition}[The graph of an orbit valuation]\label{def:enc-graphoforbits}
For $\omega \colon \mathbb{N} \to \{0,1\}$ let $G_\omega$ be the graph on
$\{0,\dots,98\}$ in which distinct $u, v$ are adjacent exactly when
$\omega(\mathrm{orb}(u,v)) = 1$.  Only the values $\omega(0),\dots,\omega(692)$
matter, so $G_\omega$ is determined by a word in $\word{693}$.
\lean{ConwayO7.graphOfOrbits}
\leanfile{ConwayO7/Orbits.lean}\leanok
\uses{def:enc-orbitmap, lem:orbits}
\end{definition}

\begin{lemma}[Invariance of the induced graph]\label{lem:enc-invariance}
For every $\omega$ the graph $G_\omega$ is $\sigma_0$-invariant:
$\sigma_0(u) \sim \sigma_0(v)$ in $G_\omega$ if and only if $u \sim v$.
\lean{ConwayO7.graphOfOrbits_invariant}
\leanfile{ConwayO7/Orbits.lean}\leanok
\uses{def:enc-graphoforbits, def:sigma0}
\end{lemma}

\begin{proof}\leanok
Immediate from the constancy of $\mathrm{orb}$ along the action
(Definition~\ref{def:enc-orbitmap}) and injectivity of $\sigma_0$.
\end{proof}

\begin{lemma}[Faithfulness of the orbit reduction]\label{lem:enc-faithful}
Let $H$ be a $\sigma_0$-invariant graph on $\{0,\dots,98\}$ and define
$\omega_H(o) = 1$ exactly when the representative pair $\mathrm{rep}(o)$ is an
edge of $H$.  Then $G_{\omega_H} = H$: an invariant graph is reconstructed
exactly from its $693$ orbit bits.
\lean{ConwayO7.orbitFunOf, ConwayO7.graphOfOrbits_orbitFunOf}
\leanfile{ConwayO7/Orbits.lean}\leanok
\uses{def:enc-graphoforbits, lem:orbits}
\end{lemma}

\begin{proof}\leanok
Given distinct $u,v$, the pair $\{u,v\}$ is reached from
$\mathrm{rep}(\mathrm{orb}(u,v))$ by iterating the action
(Lemma~\ref{lem:orbits}); each step preserves adjacency in $H$ by invariance,
so $u \sim_H v$ holds if and only if the representative pair is an edge, which
is the definition of $\omega_H(\mathrm{orb}(u,v))$.
\end{proof}

\begin{proposition}[Invariant graphs are orbit functions]\label{prop:orbitfun}
A graph $H$ on $\{0,\dots,98\}$ is $\sigma_0$-invariant if and only if
$H = G_\omega$ for some $\omega \colon \mathbb{N} \to \{0,1\}$.  This is the
verified reduction from $4851$ edge variables to $693$ orbit variables.
\lean{ConwayO7.invariant_iff_exists_orbitFun}
\leanfile{ConwayO7/Orbits.lean}\leanok
\uses{lem:enc-invariance, lem:enc-faithful}
\end{proposition}

\begin{proof}\leanok
Combine Lemma~\ref{lem:enc-invariance} (backward direction) with
Lemma~\ref{lem:enc-faithful} (forward direction, witness $\omega_H$).
\end{proof}

\section{The formula}\label{ssec:enc-formula}

\begin{definition}[Emission states]\label{def:enc-state}
An \emph{emission state} is a pair $s = (t, C)$ of a threshold $t \in \mathbb{N}$
--- the highest variable index allocated so far --- and a finite sequence $C$
of clauses, each clause a finite list of nonzero integer literals.  A
valuation $a \colon \mathbb{N} \to \{0,1\}$ \emph{satisfies} $s$ if it satisfies
every clause of $C$ (a positive literal $v$ is true when $a(v) = 1$, a
negative literal $-v$ when $a(v) = 0$).  The state is \emph{well-formed} if
$t \geq 693$ and every literal of every clause mentions a variable $\leq t$.
The generator below starts from the initial state $(693, \varnothing)$, whose
first $693$ variables $x_1,\dots,x_{693}$ are the orbit variables, and only
ever appends clauses and raises the threshold.
\lean{ConwayO7.Encoder.St, ConwayO7.Encoder.SatAll, ConwayO7.Encoder.WF}
\leanfile{ConwayO7/Encoder.lean}\leanok
\end{definition}

\begin{definition}[Orbit literals and orbit valuations]\label{def:enc-evar}
For distinct vertices $u,v$ the \emph{orbit literal} of the pair is the
(positive) literal of variable $\mathrm{orb}(u,v) + 1$; thus the adjacency of
$u,v$ in $G_\omega$ is expressed by one of the variables
$x_1,\dots,x_{693}$.  Conversely each $\omega \colon \mathbb{N} \to \{0,1\}$ induces
the valuation $a_\omega(v) = \omega(v-1)$, under which the orbit literal of
$(u,v)$ is true exactly when $u \sim v$ in $G_\omega$.
\lean{ConwayO7.Encoder.evarI, ConwayO7.Encoder.assignOfOrbits,
ConwayO7.Encoder.litSat_evarI}
\leanfile{ConwayO7/Encoder.lean}\leanok
\uses{def:enc-orbitmap, def:enc-graphoforbits, def:enc-state}
\end{definition}

\begin{definition}[The formula $\Phi$]\label{def:phi}
$\Phi$ is a propositional formula in conjunctive normal form over the orbit
variables $x_1,\dots,x_{693}$ together with auxiliary counter variables,
obtained by running four clause-emitting sections from the initial state and
reading off the final clause sequence.  Its clauses express, for the graph
$G_\omega$ associated to a valuation $\omega$ of $x_1,\dots,x_{693}$:
\begin{enumerate}
\item \emph{regularity}: every vertex has exactly $14$ neighbours;
\item \emph{common-neighbour counts}: for each orbit representative pair
$(u,v)$, the number of common neighbours of $u,v$ equals $2$ minus the
adjacency bit of $(u,v)$ (i.e.\ $\lambda = 1$, $\mu = 2$);
\item \emph{normalisation}: the neighbourhood of the fixed vertex $0$ is the
union of the first two $7$-blocks;
\item \emph{minimality}: the word $\omega \in \word{693}$ is lexicographically
minimal in its orbit under the $101$ relabelling generators of
Definition~\ref{def:relabel}.
\end{enumerate}
The cardinality conditions are rendered as clauses through standard counting
networks (a sequential counter and a modulo-$9$ totalizer, \S\ref{ssec:enc-networks});
the minimality conditions through chained lexicographic comparisons.  $\Phi$
has $432{,}882$ variables and $1{,}131{,}708$ clauses.
\lean{ConwayO7.Encoder.mkO7St, ConwayO7.Encoder.mkO7Cnf}
\leanfile{ConwayO7/Encoder.lean}\leanok
\uses{def:enc-state, def:enc-evar, lem:orbits, def:relabel, lem:sym1}
\end{definition}

\begin{lemma}[No degenerate literals]\label{lem:enc-nozero}
No clause of $\Phi$ contains the literal $0$; every literal denotes a genuine
variable with a polarity.
\lean{ConwayO7.Encoder.mkO7_no_zero}
\leanfile{ConwayO7/SatFrame.lean}\leanok
\uses{def:phi}
\end{lemma}

\begin{proof}\leanok
Checked by evaluating the generator and scanning all emitted literals.
\end{proof}

\begin{lemma}[Satisfaction transfers to the formula]\label{lem:enc-cnfsat}
A valuation satisfies the final emission state of the generator (in the
integer-literal sense of Definition~\ref{def:enc-state}) if and only if the
corresponding $0$-based valuation satisfies $\Phi$ as a formal CNF.  In
particular, if some valuation satisfies every generated clause then $\Phi$ is
satisfiable.
\lean{ConwayO7.Encoder.satAll_iff_cnfSat,
ConwayO7.Encoder.sat_mkO7Cnf_of_satAll}
\leanfile{ConwayO7/SatFrame.lean}\leanok
\uses{def:phi, lem:enc-nozero}
\end{lemma}

\begin{proof}\leanok
The two satisfaction relations agree literal by literal once no literal is
$0$ (Lemma~\ref{lem:enc-nozero}); the translation is the shift by one between
$1$-based and $0$-based variable indexing.
\end{proof}

\begin{fact}[Byte identity]\label{fact:bytes}
The formula generated inside Lean equals the parsed content of the
distributed file \texttt{o7.cnf} (SHA-256 pinned), and re-serialising that
content reproduces the file byte for byte, header included.  All certificates
of later chapters therefore speak about exactly the formula $\Phi$ of
Definition~\ref{def:phi}.
\lean{ConwayO7.mkO7Cnf_eq, ConwayO7.o7cnf_bytes}
\leanfile{ConwayO7/Data/EncoderMatch.lean}\leanok
\uses{def:phi}
\end{fact}

\section{The verification framework}\label{ssec:enc-framework}

The generator is a composition of dozens of clause-emitting transformations.
Their correctness is captured by a single mathematical notion.

\begin{definition}[Emission triples]\label{def:enc-gtrip}
Let $P$ and $Q$ be conditions on valuations and let $f$ be a transformation
of emission states.  The \emph{emission triple} $\{P\}\,f\,\{Q\}$ holds if
$f$ never lowers the threshold and, for every well-formed state $s$ with
threshold $t$, every valuation $a$ satisfying $s$ and the precondition $P$,
there exists a valuation $a'$ such that:
\begin{enumerate}
\item $a'(v) = a(v)$ for all $v \leq t$ (only \emph{fresh} auxiliary
variables, above the threshold, are reassigned);
\item $a'$ satisfies every clause of $f(s)$ --- the old clauses and the newly
emitted ones;
\item $f(s)$ is again well-formed, and $Q(a')$ holds.
\end{enumerate}
Thus a triple asserts: whatever clauses have been emitted so far, the new
gadget's clauses can always be satisfied by choosing values for its own fresh
auxiliary variables, provided the semantic condition $P$ holds on the
protected variables.
\lean{ConwayO7.Encoder.GTrip}
\leanfile{ConwayO7/SatFrame.lean}\leanok
\uses{def:enc-state}
\end{definition}

\begin{lemma}[Composition of emission triples]\label{lem:enc-gtrip-comp}
Emission triples compose: from $\{P\}\,f\,\{Q\}$ and $\{Q\}\,g\,\{R\}$
follows $\{P\}\,g \circ f\,\{R\}$.  More generally, if a condition
$\mathrm{Inv}$ is an invariant of every step of a finite iteration, it is an
invariant of the whole iteration.
\lean{ConwayO7.Encoder.GTrip.comp, ConwayO7.Encoder.GTrip.foldl_inv}
\leanfile{ConwayO7/SatFrame.lean}\leanok
\uses{def:enc-gtrip}
\end{lemma}

\begin{proof}\leanok
Chain the two witnesses; the agreement below the original threshold persists
because $f$ only raises it.  The iteration statement is induction on the list
of steps.
\end{proof}

The three sections preceding the minimality constraints have the following
semantic preconditions, stated purely as counting conditions on a valuation.

\begin{definition}[Degree invariant]\label{def:enc-deginv}
A valuation $a$ satisfies the \emph{degree invariant} if for every vertex
$u$ exactly $14$ of the $98$ orbit literals
$\{\,\mathrm{orb}(u,v)+1 : v \neq u\,\}$ are true under $a$; i.e.\ every
vertex of the graph read off the orbit variables has degree $14$.
\lean{ConwayO7.Encoder.DegInv}
\leanfile{ConwayO7/DegSpec.lean}\leanok
\uses{def:enc-evar}
\end{definition}

\begin{definition}[Common-neighbour invariant]\label{def:enc-cninv}
For a representative pair $r = (r_1,r_2)$ and a valuation $a$, let $c_r(a)$
be the number of vertices $w \notin \{r_1,r_2\}$ whose orbit literals with
both $r_1$ and $r_2$ are true, \emph{plus} $1$ if the orbit literal of
$(r_1,r_2)$ itself is true.  A valuation satisfies the
\emph{common-neighbour invariant} if $c_{\mathrm{rep}(o)}(a) = 2$ for every
orbit $o < 693$.  (Since $\mu - \lambda = 1$ for the parameters
$(99,14,1,2)$, the adjacent case $\lambda = 1$ and the non-adjacent case
$\mu = 2$ are captured by the single value $2$.)
\lean{ConwayO7.Encoder.cnCountR, ConwayO7.Encoder.CnInv}
\leanfile{ConwayO7/CnSpec.lean}\leanok
\uses{def:enc-evar, def:enc-orbitmap}
\end{definition}

\begin{definition}[Level-1 symmetry condition]\label{def:enc-sym1holds}
A valuation $a$ satisfies the \emph{level-1 symmetry condition} if for each
$i < 14$ the orbit variable of the pair $(0,\, 1 + 7i)$ has value $1$ exactly
when $i < 2$: the fixed vertex is adjacent to precisely the first two
$7$-cycles.  This is the semantic content of the normalisation of
Lemma~\ref{lem:sym1}.
\lean{ConwayO7.Encoder.Sym1Holds}
\leanfile{ConwayO7/SatFrame.lean}\leanok
\uses{def:enc-evar, lem:sym1}
\end{definition}

\begin{definition}[The pre-lex invariant]\label{def:enc-prelexinv}
The \emph{pre-lex invariant} is the conjunction of the degree invariant, the
common-neighbour invariant, and the level-1 symmetry condition.
\lean{ConwayO7.Encoder.PreLexInv}
\leanfile{ConwayO7/SectionsCompose.lean}\leanok
\uses{def:enc-deginv, def:enc-cninv, def:enc-sym1holds}
\end{definition}

\section{Soundness of the counting networks}\label{ssec:enc-networks}

The cardinality constraints of $\Phi$ are built from three networks.  Their
soundness statements all have the same shape: \emph{whenever the semantic
count is within the stated bound, the emitted clauses admit an extension of
the valuation assigning only fresh auxiliary variables}.

\begin{definition}[Unary representation of a count]\label{def:enc-unary}
A vector of literals $(y_1,\dots,y_m)$ \emph{carries the unary
representation} of a number $k$ under a valuation $a$ if $y_i$ is true
exactly when $i \leq k$; when $k$ is the number of true literals among a
vector of inputs, we say the register carries the unary count of those
inputs.
\lean{ConwayO7.TotSpec.UnaryOut, ConwayO7.MtoSpec.UnaryOf}
\leanfile{ConwayO7/TotSpec.lean}\leanok
\uses{def:enc-state}
\end{definition}

\begin{proposition}[Soundness of the sequential counter]\label{prop:enc-seqcounter}
Fix input literals $\ell_1,\dots,\ell_n$ over variables $\leq T$ and a
threshold $1 \leq \tau$ with $\tau + 2 \leq n$.  For any valuation $a$ under
which at most $\tau$ of the inputs are true, the valuation extending $a$
that assigns to each fresh register variable $s_{k,j}$ the truth value of
``more than $k$ of the first $k + j + 1$ inputs are true'' satisfies every
clause of the sequential at-most-$\tau$ counter, and agrees with $a$ on all
variables $\leq T$.
\lean{ConwayO7.SeqSpec.spec_sat}
\leanfile{ConwayO7/SeqSpec.lean}\leanok
\uses{def:enc-state}
\end{proposition}

\begin{proof}\leanok
Each clause relates a register to its predecessor and one input, and states
one instance of the recurrence for the running count
$m \mapsto \#\{i \leq m : a(\ell_i) = 1\}$; the chosen witness makes every
such instance a true arithmetic fact, the final clauses being exactly the
bound ``the count never exceeds $\tau$''.
\end{proof}

\begin{lemma}[The totalizer merge step]\label{lem:enc-totmerge}
Suppose, under a valuation $a$, two child registers carry the unary counts of
their respective input sets and the parent register carries the unary count
of the union.  Then all merge clauses of the internal node --- which say that
$i$ true literals on one side and $j$ on the other force at least $i + j$ in
the parent, and conversely --- are satisfied by $a$ itself; no fresh
variables are involved in this step.
\lean{ConwayO7.TotSpec.toUA_sat}
\leanfile{ConwayO7/TotSpec.lean}\leanok
\uses{def:enc-unary}
\end{lemma}

\begin{proposition}[Soundness of the totalizer]\label{prop:enc-totalizer}
For any well-formed state $s$ with threshold $t$, valuation $a$ satisfying
$s$, and input literals over variables $\leq t$: there is a valuation $a'$
agreeing with $a$ on all variables $\leq t$ that satisfies every clause of
the totalizer built on those inputs, and under which the output register
carries the unary count of the inputs.
\lean{ConwayO7.TotSpec.toTO_sat}
\leanfile{ConwayO7/TotSpec.lean}\leanok
\uses{def:enc-unary, lem:enc-totmerge}
\end{proposition}

\begin{proof}\leanok
Structural induction on the balanced binary tree over the inputs.  At each
internal node the fresh output register is assigned the unary representation
of the number of true leaves below it; the count at a node is the sum of the
counts at its children, so Lemma~\ref{lem:enc-totmerge} discharges the merge
clauses, and the induction hypothesis handles the subtrees.
\end{proof}

\begin{proposition}[Soundness of the modulo-9 totalizer]\label{prop:enc-mto}
Let $C$ be the number of true inputs under $a$ (at least $9$ inputs, over
variables $\leq t$).  There is an extension $a'$ of $a$ by fresh variables
satisfying all clauses of the modulo-$9$ totalizer, under which the two
output registers carry the unary representations of the quotient
$\lfloor C/9 \rfloor$ and the remainder $C \bmod 9$ respectively.
\lean{ConwayO7.MtoSpec.mtoMTO_sat}
\leanfile{ConwayO7/MtoSpec.lean}\leanok
\uses{def:enc-unary, prop:enc-totalizer}
\end{proposition}

\begin{proof}\leanok
Again induction over a binary tree, but each internal node now merges two
quotient/remainder pairs: the remainders are added as unary numbers, a carry
bit records whether their sum reached $9$, and the quotients are added
together with the carry.  The witness assigns every fresh register its exact
semantic value, so each merge clause becomes an instance of the identity
$C = 9\lfloor C/9\rfloor + (C \bmod 9)$ for the subtree counts.
\end{proof}

\begin{proposition}[The at-most-$k$ constraint via the modulo totalizer]\label{prop:enc-mto-atmost}
Fix $98$ input literals over variables $\leq t$ and $0 < k < 98$.  If at most
$k$ of the inputs are true under $a$, then the modulo-$9$ totalizer together
with its comparator clauses --- which forbid every quotient/remainder
combination exceeding $k$ --- is satisfied by an extension of $a$ assigning
only fresh variables.
\lean{ConwayO7.MtoSpec.mtoAtMost_sat}
\leanfile{ConwayO7/MtoSpec.lean}\leanok
\uses{prop:enc-mto}
\end{proposition}

\begin{proof}\leanok
Take the extension of Proposition~\ref{prop:enc-mto}.  A comparator clause
rules out the joint configuration ``quotient $\geq q$ and remainder
$\geq r$'' for each pair with $9q + r > k$; since the registers carry the
true quotient and remainder of a count $C \leq k$, no forbidden
configuration occurs.
\end{proof}

\begin{lemma}[The degree gadget]\label{lem:enc-deg-gadget}
Fix a vertex $u$ and a valuation $a$ under which exactly $14$ of the $98$
orbit literals at $u$ are true.  Then the exactly-$14$ gadget for $u$ --- an
at-least-$14$ network run on the negated literals (as at-most-$84$) followed
by an at-most-$14$ network, both modulo-$9$ totalizers --- admits an
extension of $a$ by fresh variables satisfying all its clauses.
\lean{ConwayO7.Encoder.encEqualsDeg_sat}
\leanfile{ConwayO7/DegSpec.lean}\leanok
\uses{prop:enc-mto-atmost, def:enc-deginv, def:enc-evar}
\end{lemma}

\begin{proof}\leanok
Exactly $14$ of $98$ literals true means exactly $84$ of the negated
literals are true; apply Proposition~\ref{prop:enc-mto-atmost} twice in
sequence, composing the two fresh extensions.
\end{proof}

\begin{proposition}[The degree section]\label{prop:enc-deg-section}
The degree invariant is a valid pre- and postcondition for the degree gadget
of each vertex, and hence for the entire degree section --- the iteration of
the gadget over all $99$ vertices --- as one emission triple.
\lean{ConwayO7.Encoder.degreeStep_trip, ConwayO7.Encoder.deg_section_trip}
\leanfile{ConwayO7/DegSpec.lean}\leanok
\uses{def:enc-gtrip, lem:enc-gtrip-comp, lem:enc-deg-gadget}
\end{proposition}

\begin{proof}\leanok
The degree invariant concerns only the orbit variables ($\leq 693$), which
every gadget's extension leaves untouched; so it survives each step, and
Lemma~\ref{lem:enc-deg-gadget} supplies each step's satisfiability.  Compose
by Lemma~\ref{lem:enc-gtrip-comp}.
\end{proof}

\begin{lemma}[The common-neighbour products]\label{lem:enc-cn-product}
Fix a representative pair $r = (r_1,r_2)$.  For each candidate common
neighbour $w$ the gadget allocates a fresh variable $p_w$ with three clauses
expressing $p_w \leftrightarrow (e_{r_1 w} \wedge e_{r_2 w})$ in terms of the
orbit literals.  Assigning each $p_w$ its semantic value --- the conjunction
of the two adjacencies --- satisfies all product clauses, and the recorded
list of product literals evaluates to exactly the indicator sequence of
common neighbourhood.
\lean{ConwayO7.Encoder.prodFold_sat}
\leanfile{ConwayO7/CnSpec.lean}\leanok
\uses{def:enc-evar, def:enc-cninv}
\end{lemma}

\begin{lemma}[The exactly-2 gadget]\label{lem:enc-cn-equals}
Fix $98$ literals over variables $\leq t$ of which exactly two are true
under $a$.  Then the exactly-$2$ gadget --- a sequential at-most-$96$
counter on the negated literals (enforcing at least $2$) followed by a
sequential at-most-$2$ counter --- admits an extension of $a$ by fresh
variables satisfying all its clauses.
\lean{ConwayO7.Encoder.specSeqEquals_sat}
\leanfile{ConwayO7/CnSpec.lean}\leanok
\uses{prop:enc-seqcounter}
\end{lemma}

\begin{proof}\leanok
Exactly $2$ true literals means exactly $96$ true negations; apply
Proposition~\ref{prop:enc-seqcounter} to each counter in turn and compose
the extensions.
\end{proof}

\begin{proposition}[The common-neighbour section]\label{prop:enc-cn-section}
The common-neighbour invariant is a valid pre- and postcondition for the
gadget of each orbit --- the product stage followed by the exactly-$2$
gadget on the $97$ product literals together with the pair's own orbit
literal --- and hence for the whole section over all $693$ orbits, as one
emission triple.
\lean{ConwayO7.Encoder.specCnStep_trip, ConwayO7.Encoder.cn_section_trip}
\leanfile{ConwayO7/CnSpec.lean}\leanok
\uses{def:enc-gtrip, lem:enc-gtrip-comp, lem:enc-cn-product,
lem:enc-cn-equals, def:enc-cninv}
\end{proposition}

\begin{proof}\leanok
By Lemma~\ref{lem:enc-cn-product} the product literals plus the adjacency
literal have exactly $c_r(a)$ true entries, which the invariant fixes at
$2$; Lemma~\ref{lem:enc-cn-equals} then extends the valuation over the two
counters.  The invariant mentions only orbit variables and survives every
fresh extension; compose by Lemma~\ref{lem:enc-gtrip-comp}.
\end{proof}

\begin{lemma}[The normalisation section]\label{lem:enc-sym1-section}
The level-1 symmetry condition is a valid pre- and postcondition for the
section emitting the $14$ unit clauses that pin the neighbourhood of the
fixed vertex: under the condition, each unit clause is already true and no
fresh variables are needed.
\lean{ConwayO7.Encoder.sym1_section_trip}
\leanfile{ConwayO7/SatFrame.lean}\leanok
\uses{def:enc-gtrip, lem:enc-gtrip-comp, def:enc-sym1holds}
\end{lemma}

\begin{proposition}[Pre-lex satisfiability]\label{prop:counters}
For every valuation of the orbit variables satisfying the pre-lex invariant
(Definition~\ref{def:enc-prelexinv}) --- that is, the arithmetic conditions
(1)--(3) of Definition~\ref{def:phi} --- there is an extension assigning
only auxiliary variables above $693$ that satisfies every clause emitted by
the degree, common-neighbour, and normalisation sections of the generator.
\lean{ConwayO7.Encoder.preLexSpec_trip, ConwayO7.Encoder.preLex_sat}
\leanfile{ConwayO7/SectionsCompose.lean}\leanok
\uses{def:enc-prelexinv, prop:enc-deg-section, prop:enc-cn-section,
lem:enc-sym1-section, def:phi}
\end{proposition}

\begin{proof}\leanok
Each of the three component invariants is stable under fresh extension, so
each section's triple upgrades to a triple for the full pre-lex invariant;
compose the three (Lemma~\ref{lem:enc-gtrip-comp}) and instantiate at the
initial state $(693, \varnothing)$.  A separate computation identifies the
state built from the specification-level common-neighbour gadgets with the
state the concrete generator emits, so the conclusion holds for the clauses
of $\Phi$ itself.
\end{proof}

\section{The bridge from strong regularity}\label{ssec:enc-bridge}

Finally, the semantic preconditions above are not ad hoc: they are exactly
what strong regularity delivers.

\begin{lemma}[Regularity yields the degree invariant]\label{lem:enc-deg-bridge}
If $G_\omega$ is strongly regular with parameters $(99,14,1,2)$, then the
induced valuation $a_\omega$ satisfies the degree invariant: for each vertex
$u$ the number of true orbit literals at $u$ is the degree of $u$ in
$G_\omega$, which is $14$.
\lean{ConwayO7.Encoder.degInv_of_srg}
\leanfile{ConwayO7/Bridge.lean}\leanok
\uses{def:srg, def:enc-graphoforbits, def:enc-evar, def:enc-deginv}
\end{lemma}

\begin{lemma}[The srg identities yield the common-neighbour invariant]\label{lem:enc-cn-bridge}
If $G_\omega$ is strongly regular with parameters $(99,14,1,2)$, then
$a_\omega$ satisfies the common-neighbour invariant: for a representative
pair $r$, the count $c_r(a_\omega)$ equals the number of common neighbours
of $r_1, r_2$ plus the adjacency bit, i.e.\ $\lambda + 1 = 2$ when the pair
is an edge and $\mu = 2$ when it is not.
\lean{ConwayO7.Encoder.cnInv_of_srg}
\leanfile{ConwayO7/Bridge.lean}\leanok
\uses{def:srg, def:enc-graphoforbits, def:enc-evar, def:enc-cninv}
\end{lemma}

\begin{proof}\leanok
The candidate list ranges over all vertices other than $r_1, r_2$, and a
candidate's two orbit literals are true exactly when it is a common
neighbour in $G_\omega$; the count is therefore
$|N(r_1) \cap N(r_2)|$ plus the adjacency indicator, and both cases of the
strong-regularity identity give the value $2$ since $\mu - \lambda = 1$.
\end{proof}

\begin{proposition}[Strong regularity yields the pre-lex invariant]\label{prop:enc-prelex-bridge}
If $G_\omega$ is strongly regular with parameters $(99,14,1,2)$ and
$a_\omega$ satisfies the level-1 symmetry condition, then $a_\omega$
satisfies the full pre-lex invariant.  Combined with
Proposition~\ref{prop:counters}, every normalised $\sigma_0$-invariant srg
would yield a valuation satisfying all pre-normalisation clauses of $\Phi$;
the minimality clauses are handled in Chapter~\ref{sec:lex}.
\lean{ConwayO7.Encoder.preLexInv_of_srg}
\leanfile{ConwayO7/Bridge.lean}\leanok
\uses{lem:enc-deg-bridge, lem:enc-cn-bridge, def:enc-sym1holds,
def:enc-prelexinv}
\end{proposition}

\begin{proof}\leanok
Conjoin Lemmas~\ref{lem:enc-deg-bridge} and~\ref{lem:enc-cn-bridge} with the
assumed symmetry condition.
\end{proof}

%=====================================================================
